\section{Held-out Routing Interventions}
\label{sec:guiding}

\paragraph{Calibration and evaluation population.}
Both policies split source problems within each benchmark, allocating 1,082
problems to calibration and 465 to evaluation (70/30, seed 42).
Calibration uses 3,014 scored traces from 1,001 problems; 238 traces with
missing saved correctness are excluded. This differs from the correlation
analyzer, which counts known-gold missing answers as incorrect. The
problem-weighted calibration accuracy is 0.5817 on this scored subset.
Consequently, calibration targets may be sensitive to answer-extraction
failures and must not be interpreted as estimates over all allocated traces.
Each problem contributes total weight one, divided across its usable attempts.
Evaluation retains 150 MATH-500, 9 AIME24, 9 AIME25, 203 OlympiadBench,
12 AMC23, and 82 Minerva problems. With 32 attempts on each AIME/AMC
problem, this gives 1,395 paired attempts. No evaluation problem is used
in policy calibration.

\paragraph{Identity policy.}
At each of the six selected Qwen layers, expert inclusion is weighted by
its routing fraction and problem weight. Experts require support from at
least four distinct problems and positive accuracy lift over the calibration
baseline. At most eight identities per layer are selected; positive lifts
are normalized to the largest retained lift. A signed strength scales the
additive router-logit bias. Strength $+1$ promotes these identities;
strength $-1$ suppresses the same identities. The negative control does
not learn a separate set of incorrectly associated experts. Unlike the
pooled-ID observational screen, policy targets are layer-qualified.

\paragraph{Margin policy.}
The policy calibrates the trace-mean selection-boundary margin
$p_{(k)}-p_{(k+1)}$, with $k=8$, using five problem-weighted quantile bins.
Among bins with at least four supporting problems, it selects the
highest-accuracy positive-lift bin, breaking ties by support and then lower
margin. At inference, probability mixing moves the token-level gap toward
the selected interval: uniform mixing reduces the gap, whereas mixing with
a distribution uniform over the selected top-$k$ increases it.
Trace-level calibration and token-level application have different units.
Both policies transfer routing targets from Unsloth 4-bit replay to GPTQ
generation, which is an additional unisolated deployment difference.

\paragraph{Matched evaluation.}
Each policy is compared with its own baseline using identical problem IDs,
prompts, prompt token IDs, per-attempt seeds, sampling settings, engine
configuration, and recorded software versions. Generation uses temperature
0.6, top-$p$ 0.95, and a 32,768-token limit. The identity $+1$ and margin
comparisons use maximum concurrency 16; the identity $-1$ comparison uses
25. Baseline outputs therefore cannot be treated as interchangeable across
comparisons. Only complete matched stochastic pairs are reported; the
strength-$+2$ condition has no complete matched evaluation artifact.
Separate greedy and restricted-benchmark outputs are not pooled here.
The engine writes completed records incrementally and can resume without
continuing a partially generated sequence. Baseline reuse requires matching
contracts and a verified output hash.

\begin{table*}[ht]
\centering\small
\caption{Complete matched Qwen3.5 GPTQ stochastic evaluations, 1,395 attempts on 465 held-out problems per comparison. B/G denotes baseline/guided; changes are percentage points. Flips are wrong-to-right/right-to-wrong. These are descriptive point estimates without cluster confidence intervals.}
\label{tab:guiding-summary}
\begin{tabular}{lrrrrrr}
\toprule
Policy & B accuracy & G accuracy & Change & Flips & Limit hits B/G & Mean tokens B/G \\
\midrule
Identity $+1$ & 0.7391 & 0.7355 & -0.36 & 39/44 & 140/146 & 13863/13706 \\
Identity $-1$ & 0.7319 & 0.7398 & +0.79 & 47/36 & 154/138 & 13830/13919 \\
Margin $1$ & 0.7384 & 0.7369 & -0.14 & 39/41 & 144/143 & 13848/13840 \\
\bottomrule
\end{tabular}
\end{table*}

\begin{table*}[ht]
\centering\small
\caption{Qwen3.5 GPTQ guiding results by benchmark. Correct counts are baseline/guided; accuracy changes are percentage points. Repeated attempts from one source problem are dependent.}
\label{tab:guiding-datasets}
\begin{tabular}{llrrr}
\toprule
Policy & Dataset & Attempts & Correct B/G & Change \\
\midrule
Identity $+1$ & aime24 & 288 & 185/177 & -2.78 \\
Identity $+1$ & aime25 & 288 & 246/246 & +0.00 \\
Identity $+1$ & amc23 & 384 & 381/381 & +0.00 \\
Identity $+1$ & math500 & 150 & 106/106 & +0.00 \\
Identity $+1$ & minerva & 82 & 16/18 & +2.44 \\
Identity $+1$ & olympiad & 203 & 97/98 & +0.49 \\
Identity $-1$ & aime24 & 288 & 178/182 & +1.39 \\
Identity $-1$ & aime25 & 288 & 242/248 & +2.08 \\
Identity $-1$ & amc23 & 384 & 382/382 & +0.00 \\
Identity $-1$ & math500 & 150 & 103/102 & -0.67 \\
Identity $-1$ & minerva & 82 & 20/19 & -1.22 \\
Identity $-1$ & olympiad & 203 & 96/99 & +1.48 \\
Margin $1$ & aime24 & 288 & 180/184 & +1.39 \\
Margin $1$ & aime25 & 288 & 243/246 & +1.04 \\
Margin $1$ & amc23 & 384 & 380/380 & +0.00 \\
Margin $1$ & math500 & 150 & 106/104 & -1.33 \\
Margin $1$ & minerva & 82 & 20/18 & -2.44 \\
Margin $1$ & olympiad & 203 & 101/96 & -2.46 \\
\bottomrule
\end{tabular}
\end{table*}


\paragraph{Observed effects and limits.}
Identity $+1$ changes 83 correctness outcomes and loses five net correct
answers; margin guiding changes 80 and loses two. Identity $-1$ changes 83
and gains eleven. The corresponding pooled changes are $-0.36$, $-0.14$,
and $+0.79$ percentage points. This does not establish a reliable benefit
from promoting positive-lift experts or targeting the calibrated margin.
The negative-strength point estimate is positive, but differences across
strengths are not a controlled ranking because their baselines and
concurrency differ. These are exploratory comparisons without paired
source-problem confidence intervals or a separate strength-selection set.
Pooled attempt accuracy also weights each AIME/AMC problem 32 times as
heavily as each single-attempt problem; benchmark counts are reported to
make this weighting explicit.

\paragraph{What the margin hook actually changes.}
Full diagnostics for the completed margin comparison record 20,657,039
hook token evaluations per selected layer. The top-eight identity set changes
on 32.7--41.8\% of these evaluations. Mean distance to the target interval
falls from $2.19$--$4.71\times10^{-4}$ to
$0.285$--$0.902\times10^{-4}$ across layers. Thus the numerical operation
moves the intended gap, but it does not isolate a membership-preserving
intervention after finite-precision log-probability conversion and casting.
These diagnostics describe the hooked logits, not an independent recording
of the kernel's dispatched expert IDs. Their denominator can include
prefill, padding, and recomputation and is not a generated-token count.
The identity $-1$ comparison uses minimal diagnostics, so detailed hook
statistics are not available on a common basis across all conditions.
